% Configurazione
\documentclass{article}

\usepackage{titling} % Required for inserting the subtitle
\usepackage{graphicx} % Required for inserting images
\usepackage{tabularx} % Per l'ambiente tabularx (tabelle)
\usepackage{calc} % Sempre per le tabelle
\usepackage[hidelinks]{hyperref} % Per i collegamenti ipertestuali, ad esempio sulla table of contents
\usepackage{xcolor} % Per colorare il testo
\usepackage{colortbl} % Per colorare le celle delle tabelle
\usepackage{lipsum} % Per generare lorem ipsum
\usepackage[normalem]{ulem} % Per sottolineare il testo
\usepackage{array} % Per la visualizzazione fluttuante di array di domande e risposte
\usepackage{ragged2e} % Pacchetto necessario per \justifying che giustifica il testo di tabelle

\newcommand{\ulhref}[2]{\href{#1}{\uline{#2}}} % Nuovo comando per sottolineare i link
\newcommand{\ulref}[1]{\uline{\ref{#1}}} % Nuovo comando per sottolineare i collegamenti a immagini
\setlength{\parindent}{0pt} % Rimuove il rientro automatico dei paragrafi

\graphicspath{ {immagini/} {../../../../shared/images/} }


% Variabili
\newcommand{\data}{15 Ottobre 2025}


% Struttura
\begin{document}

% Parte costante
\begin{minipage}{0.4\textwidth}
    \includegraphics[width=0.6\textwidth]{logo_unipd.png}
\end{minipage}
\begin{minipage}{0.55\textwidth}
    \textcolor{red}{\textbf{Università degli Studi di Padova}} \\
    \textcolor{red}{Laurea: Informatica} \\
    \textcolor{red}{Corso: Ingegneria del Software} \\
    \textcolor{red}{Anno Accademico: 2025/2026}
\end{minipage}

\begin{minipage}{0.4\textwidth}
    \includegraphics[width=0.6\textwidth]{logo_gruppo.jpeg}
\end{minipage}
\begin{minipage}{0.55\textwidth}
    \textbf{Gruppo: NullPointers Group} \\
    Email: \textsf{groupnullpointers@gmail.com}
\end{minipage}

\vspace{2cm}
{
\centering
\Huge\bfseries Verbale Riunione\par
\vspace{0.5cm}
\Large\bfseries \data\par
}
\newpage
\tableofcontents
\newpage

% Parte variabile
\section{Registro delle modifiche}

\begin{table}[htbp]
	\begin{tabular}{|c|c|c|c|p{0.5\textwidth}|}
		\hline
		\rowcolor[gray]{0.9}
		Vers & Data & Autore & Ruolo & Descrizione \\
		\hline
		1.0 & 17-10 & Lisa Casagrande & - & Aggiunte informazioni inizio verbale + approvazione verbale\\
		\hline
		0.2 & 17-10 & Laura Pieripolli & - & Correzioni minori \\ 
		\hline
		0.12 & 16-10 & Matteo Mazzaretto & - & Modifica colore tabella e modalità di numerazione codici\\
		\hline
		0.11 & 16-10 & Luca Marcuzzo & - & Modifica errata corrige su data prossima riunione \\
		\hline
		0.10 & 16-10 & Matteo Mazzaretto & - & Aggiunta tabella decisioni \\
		\hline
		0.01 & 15-10 & Tommaso Ceron & - & Creazione documento \\
		\hline
	\end{tabular}
\end{table}
\newpage
\section{Informazioni generali}

\begin{itemize}
    \item \textbf{Tipo di riunione:} interna
    \item \textbf{Luogo:} Meeting Discord
    \item \textbf{Data:} 2025/10/15
    \item \textbf{Ora inizio:} 15:00
    \item \textbf{Ora fine:} 17:40
    \item \textbf{Responsabile:} Matteo Mazzaretto 
    \item \textbf{Scriba:} Tommaso Ceron
    \item \textbf{Partecipanti:}
    \begin{itemize}
        \renewcommand{\labelitemii}{--}
        \item Laura Pieripolli
        \item Lisa Casagrande
        \item Luca Marcuzzo
        \item Marco Brunello
        \item Matteo Mazzaretto
        \item Tommaso Ceron
    \end{itemize}
\end{itemize}
\newpage
\section{Ordine del giorno}

\begin{itemize}
    \item Definizione way of working e creazione nome/logo del gruppo
    \item Discutere sulle nostre preferenze riguardo i capitolati di progetto
\end{itemize}
\newpage
\section{Diario della riunione}

	In questa riunione abbiamo deciso il nome del gruppo e il logo. Il nostro nome è "NullPointers Group" e il nostro logo è
	\begin{figure}[h]
		\begin{center}
			\includegraphics[width=0.3\textwidth]{logo_gruppo.jpeg}
		\end{center}
	\end{figure}\\
	Abbiamo creato il recapito mail per il nostro gruppo.\\
	Abbiamo creato la repository GitHub per la documentazione ed è stato approvato l'utilizzo di LaTeX.\\
	Inoltre, abbiamo deciso di utilizzare una tecnologia che GitHub mette a disposizione, ovvero GitHub Pages.\\
	Questa tecnologia permette di creare un sito mantentendo i file necessari nella repository.\\
	Abbiamo optato per questa tecnologia poiché permette di mantenere più ordinata tutta la nostra documentazione.\\
    Ogni componente del gruppo ha espresso le proprie preferenze riguardo ai capitolati, che sono confluite in \href{https://docs.google.com/spreadsheets/d/1QaMKRtNCo9-x2O6sPm80BQgh0mwIuRUwIfmQZ6hzgeY/edit?gid=0#gid=0}{\textcolor{blue}{questo}} documento.\\
    Infine, abbiamo deciso i codici da assegnare alle decisioni e agli obiettivi futuri durante la redazione dei verbali, sia interni che esterni.\\
    Fissata la prossima riunione, il giorno 17/10/2025 alle ore 10:00, in presenza.
\newpage
\section{Decisioni}

	\begin{table}[htbp]
		\begin{tabular}{|p{0.2\textwidth}|p{0.8\textwidth}|}
			\hline
			\rowcolor[gray]{0.9}
			Codice & Decisione \\
			\hline
			VI 1.1 & Decisione codici. Utilizzeremo i codici VI (verbale interno), VE (verbale esterno) e TD (task futura), seguiti da un identificativo numerico (x.y) che specifica il numero progressivo del verbale e il punto specifico all'interno del documento. \\
			\hline
			VI 1.2	& Decisione nome gruppo e creazione logo \\
			\hline
			VI 1.3 & Creazione recapito mail gruppo \\
			\hline
			VI 1.4 & Creazione repository documentazione \\
			\hline
			VI 1.5	& Approvazione utilizzo LaTeX per documentazione \\
			\hline
			VI 1.6 & Approvazione utilizzo GitHub Pages per sito repository \\
			\hline
			VI 1.7 & Creazione sito \\
			\hline
			VI 1.8 & Espresse preferenze per i capitolati nel documento segnalato nel diario della riunione \\
			\hline
			VI 1.9 & Fissata prossima riunione \\
			\hline
		\end{tabular}
	\end{table}
\newpage
\section{Todo}

Durante la riunione sono emersi i seguenti task da svolgere:

\vspace{0.5cm}

\begin{table}[htbp]
	\begin{tabular}{|c|c|c|}
		\hline
		\rowcolor[gray]{0.9}
		Codice & Assegnatario & Task Todo \\
		\hline
		TD1 & \it{Tutti i membri del gruppo} & Design e scrittura codice HTML/CSS del sito\\
		\hline
		TD2 & \it{Tutti i membri del gruppo} & Iniziare la stesura della domanda di candidatura \\
		\hline
		TD3 & \it{Tutti i membri del gruppo} & Creazione GitHub Actions per automatismi nella repository \\
		\hline
	\end{tabular}
\end{table}

\end{document}